\section{Experiments}

In our experiments, we (1) establish a performance baseline on \ours using \sysname, (2) analyze the performance of the three components of \sysname, (3) demonstrate the importance of in-domain training data, and (4) present promising qualitative results on verifying claims about COVID-19 using \sysname.

\subsection{Results}

Table \ref{tbl:main_results} shows the full-pipeline performance of \sysname on the \ours test set, evaluated using the abstract-level and sentence-level metrics defined in \S\ref{sec:task_definition}. The \fonecoarse value of 46.5 indicates that, for roughly half of the claim-abstract pairs, \sysname correctly identifies the \supports or \refutes label \emph{and} provides reasonable evidence to justify the decision. Given the difficulty of the task and limited in-domain training data, we consider this a promising result, while leaving plenty of room for improvement.

\input{tables/new-results.tex}

\paragraph{Oracle experiments} To examine the performance of each system component, we run the \sysname pipeline, replacing some components with ``oracles'' that always make correct predictions when given correct inputs.\footnote{For instance, the ``Label'' oracle always predicts the correct label as long as it is provided with at least one full gold rationale as input.} The oracles are only used for prediction; they do not affect the training process. The first three rows in Table~\ref{tbl:main_results} are \emph{double-oracle}. These experiments isolate the performance of a single model component. Row 1 shows that the retrieval module has an abstract recall of 68.5 F1. The precision is 100 because the oracle label predictor will always predict \notenough for false positive retrievals. Row 2 shows that, given gold abstracts, the rationale selection module identifies rationale sentences with an F1 of 75.4. Row 3 shows that, given gold rationales from gold abstracts, the label classifier has a label prediction F1 of 83.3. Since the test set is balanced and gold abstracts are provided, this is roughly comparable to a classification accuracy over the \supports and \refutes label classes.

The next three rows are \emph{single-oracle}. Row 4 shows the performance of the rationale selector and label classifier together, when given gold documents as input. This is similar to the \eraser evaluation setup.

\paragraph{Training datasets} During model development, we train the \componenttwo and \componentthree modules on four different datasets: \fever, \snopes, \ours, and \fever pretraining followed by \ours fine-tuning. The \componenttwo module is evaluated on its ability to identify rationale sentences given gold abstracts.\footnote{Our \fever-trained \componenttwo module achieves 79.9 sentence-level F1 on the \fever test set, virtually identical to the value of 79.6 reported in \citet{DeYoung2019ERASERAB}}. As in Section \ref{sec:evaluation}, a rationale sentence is only identified correctly if all other sentences from its gold rationale are also identified. The \componentthree module is evaluated on its classification accuracy given gold rationales from evidence abstracts (including evidence documents labeled \notenough). The results of these experiments are shown in Table \ref{tbl:components}. For \componenttwo, training on \ours alone produces good results, perhaps because domain-specific lexical cues are sufficient in most cases for identifying rationale sentences. For the more complex reasoning involved in \componentthree, domain adaptation was the most effective approach, training first on the large \fever dataset and then the smaller in-domain \ours training set. Based on these results, we use the \componenttwo module trained on \ours only, and the \componentthree module trained on \fever + \ours for our final end-to-end system \sysname. Additional implementation details can be found in Appendix \ref{sec:model_details}.

\begin{table}[t]
  \footnotesize
  \setlength{\tabcolsep}{0.5em}
  \centering

  \begin{tabularx}{\columnwidth}{l *{3}{c}  *{3}{c}}
    \toprule
 & \multicolumn{3}{c}{\textsc{Rational-Select.}} & \multicolumn{3}{c}{\textsc{Label-Pred.}} \\
%  & \multicolumn{3}{c}{\textsc{Selection}}  & \multicolumn{3}{c}{\textsc{Prediction}} \\
    % \cmidrule(lr){2-4} \cmidrule(lr){5-7}

    \midrule
    \textbf{Training data} & P & R & F1 & \multicolumn{3}{c}{\textsc{Acc.}} \\

    \arrayrulecolor{black!30}\midrule
    % \hdashline[0.4pt/2pt]

    \fever         &      41.5 &   57.9 &  48.4 &  \multicolumn{3}{c}{67.6} \\
    \snopes        &      42.5 &   62.3 &  50.5 &  \multicolumn{3}{c}{71.3} \\
    \ours          &      73.7 &   70.5 &  \textbf{72.1} &  \multicolumn{3}{c}{75.7} \\
    \fever + \ours &      72.4 &   67.2 &  69.7 &  \multicolumn{3}{c}{\textbf{81.9}} \\

    \arrayrulecolor{black}
    \midrule

    \textbf{Sentence encoder} & P & R & F1 & \multicolumn{3}{c}{\textsc{Acc.}} \\
    \arrayrulecolor{black!30}\midrule
    % \hdashline[0.4pt/2pt]

    \scibert      &      74.5 &   74.3 &  \textbf{74.4} &  \multicolumn{3}{c}{69.2} \\
    \bioroberta   &      75.3 &   69.9 &  72.5 &  \multicolumn{3}{c}{71.7} \\
    \robertabase  &      76.1 &   66.1 &  70.8 &  \multicolumn{3}{c}{62.9} \\
    \robertalarge &      73.7 &   70.5 &  72.1 &  \multicolumn{3}{c}{\textbf{75.7}} \\

    \arrayrulecolor{black}
    \midrule

    \textbf{Model inputs} & P & R & F1 & \multicolumn{3}{c}{\textsc{Acc.}} \\
    \arrayrulecolor{black!30}\midrule
    % \hdashline[0.4pt/2pt]
    Claim-only         &      - &   - &  - &  \multicolumn{3}{c}{44.5} \\
    Abstract-only         &   60.1 &   60.9 &  60.5 &  \multicolumn{3}{c}{53.3} \\

    \arrayrulecolor{black}
    \bottomrule
  \end{tabularx}

  \caption{Comparison of different training datasets, encoders, and model inputs for \componenttwo and \componentthree, evaluated on the \ours dev set. A claim-only model cannot select rationales.}
  \label{tbl:components}

\end{table}

%%%%%%%%%%%%%%%%%%%%%%%%%%%%%%%%%%%%%%%%%%%%%%%%%%%%%%%%%%%%%%%%%%%%%%%%%%%%%%%%

\subsection{Verifying claims about COVID-19}

We conduct exploratory experiments using our system to fact-check claims concerning COVID-19. We task a medical student to write 36 COVID-related claims. For each claim $c$, we use \sysname to predict evidence abstracts $\Ehat(c)$. The same medical student annotator assigns a label to each $(c, \Ehat(c))$ pair. A pair is labeled \emph{plausible} if at least half of the evidence abstracts in $\Ehat(c)$ are judged to have reasonable rationales and labels. It is labeled \emph{missed} if $\Ehat(c) = \emptyset$. Finally, it is labeled \emph{implausible} if the majority of the abstracts in $\Ehat(c)$ have irrelevant rationales or incorrect labels.

Table \ref{tbl:covid_example} shows two example claims, both with supporting and refuting evidence identified by \sysname. For the majority of these COVID-related claims (23 out of 36), the rationales produced by \sysname was deemed \emph{plausible} by our annotator, demonstrating that \sysname is able to successfully retrieve and classify evidence in many cases. An examination of errors reveals that the system can be confused by context, where abstracts are labeled \supports or \refutes even though the rationale sentences reference a different disease or drug from the claim. An example of this is also provided in Table \ref{tbl:covid_example}.

\input{08-discussion.tex}
